\subsection{A real Montgomery/Kummer case for the toric discrepancy: Curve25519 order-four leakage}
\label{subsec:curve25519-discrepancy-experiment}

The experiments in this subsection are not intended to validate the compressed
secp256k1 residue labels used elsewhere in the paper.  They validate a weaker
but more concrete claim: when the full toric discrepancy is retained, and in
particular when its exceptional divisor is used as the side-channel label, the
resulting label coincides with a documented leakage event in a real
Montgomery/Kummer implementation.  The historical implementation case is the
Libgcrypt Curve25519 attack of Genkin, Valenta, and Yarom, ``May the Fourth Be
With You''~\cite{genkin2017mayfourth}.  The vulnerability was assigned
CVE-2017-0379 and was mitigated in Libgcrypt 1.7.9 and 1.8.1~\cite{gnupg2017cve0379,nvd2017cve0379}.

Consider a Montgomery curve
\[
  E_A : v^2 = u^3 + A u^2 + u
\]
over a prime field.  The affine Kummer doubling map is
\[
  D_A(u)=\frac{(u^2-1)^2}{4u(u^2+Au+1)}.
\]
We compare this map with the same tangent/toric reference law used in our
model,
\[
  T(u)=\frac{2u}{1-u^2}.
\]
Using the toric subtraction law, define the discrepancy
\[
  \Delta_A(u)=D_A(u)\ominus T(u)
  =\frac{D_A(u)-T(u)}{1+D_A(u)T(u)}.
\]
A direct simplification gives
\[
  \Delta_A(u)=
  \frac{u^6+5u^4+8Au^3+11u^2-1}
       {2u(u-1)(u+1)(u^2+2Au+3)} .
\]
Thus the discrepancy has poles at
\[
  u=1 \quad\text{and}\quad u=-1 .
\]
For Curve25519, where \(A=486662\) and \(p=2^{255}-19\), the class \(u=1\) is
an order-four point on the curve, whereas \(u=-1\) is an order-four point on the
quadratic twist.  These are precisely the low-order Kummer inputs used in the
published Libgcrypt attack.  In projective Kummer coordinates \(u=X/Z\), the
same exceptional condition is
\[
  u=\pm 1 \quad\Longleftrightarrow\quad X=\pm Z
  \quad\Longleftrightarrow\quad Z^2-X^2=0 .
\]
This is the explicit bridge between the toric discrepancy and the implementation
state: the side-channel label is the occurrence of a projective representative
lying on the exceptional divisor of \(\Delta_A\).

\paragraph{Experiment 1: algebraic verification of the discrepancy pole.}
The script \texttt{curve25519\_toric\_discrepancy.py} verifies the equality
between the direct definition of \(\Delta_A\) and the simplified rational form
above over \(\mathbb{F}_{2^{255}-19}\).  It performed 10,000 random consistency
checks with no mismatch.  It also evaluated the two exceptional inputs.  The
input \(u=1\) was classified as a toric pole and as a curve point; the input
\(u=-1\) was classified as a toric pole and as a quadratic-twist point.  In both
cases Montgomery/Kummer doubling itself remains defined, but the toric
comparison map is undefined.  This distinction is important: the leakage is not
caused by a failure of the Montgomery ladder formula, but by the implementation
being driven into a special Kummer state that is exceptional for the discrepancy
map.

\paragraph{Experiment 2: deterministic reproduction of the Libgcrypt line-19 label.}
The script \texttt{order4\_line19\_poc.py} implements the projective Montgomery
step used in the attack paper's simplified description of Libgcrypt.  For a
ladder state \(Q_0=(X_0,Z_0)\), the step computes
\[
  l_{11}=(X_0+Z_0)^2-(X_0-Z_0)^2=4X_0Z_0
\]
and then performs the line-19 operation
\[
  Z_d \leftarrow l_{14}l_{11}\bmod p .
\]
Libgcrypt's variable-size MPI reduction exited early when the unreduced product
had fewer limbs than the modulus.  Following the terminology of
Genkin--Valenta--Yarom, we call the line-19 reduction \emph{short} if this early
exit occurs and \emph{long} otherwise.  The binary leakage symbol used in this
subsection is
\[
  \rho_i =
  \begin{cases}
    0, & \text{short line-19 reduction at ladder iteration } i,\\
    1, & \text{long line-19 reduction at ladder iteration } i.
  \end{cases}
\]
For an order-four input, the Montgomery ladder invariant implies that one of the
two tracked Kummer states is an order-four toric-pole state and the other is the
identity or the order-two point.  Consequently, after the initial exceptional
warm-up window, the short/long predicate is equivalent to a scalar-bit
transition:
\[
  \rho_i = k_i \oplus k_{i-1} .
\]
In our deterministic run over a generated 255-bit MPI scalar with a known top bit, the script checked 251
post-warm-up ladder rows.  It found zero mismatches against the rule
\(\rho_i=k_i\oplus k_{i-1}\).  The trace contained 120 long reductions and 135
short reductions.  Given bit \(k_3\), the suffix \(k_4,\ldots,k_{254}\) was
recovered exactly from the reduction labels.  If only the fixed X25519 top bit
is used, the first three unresolved bits form an eight-candidate prefix that can
be enumerated.

\paragraph{Experiment 3: timing realization of the binary discrepancy label.}
The script \texttt{run\_timing\_surrogate.py} compiles and executes a small C
program, \texttt{timing\_surrogate.c}, that deliberately emulates the timing
difference between a short and a long modular reduction.  The program does not
implement a live Libgcrypt exploit.  It is a controlled timing source for the
same binary label \(\rho_i\).  In the current run, one scalar multiplication
trace contained 255 timing samples.  The measured means were approximately
7,994 ns for short reductions and 325,818 ns for long reductions.  A threshold
classifier separated the two classes with 100\% label accuracy and recovered the
post-warm-up scalar suffix with 100\% accuracy.

A second script, \texttt{run\_timing\_surrogate\_repeated.py}, models the
multi-trace aggregation used in the real Flush+Reload attack.  It collected 11
timing traces for the same scalar.  Each trace was classified using an
unsupervised two-cluster threshold, and the predicted \(\rho_i\) values were
majority-voted per ladder edge.  The majority-vote classifier achieved 100\%
label accuracy and recovered the post-warm-up scalar suffix exactly.  This
experiment is included because a single bit-transition error propagates through
the recurrence \(k_i=k_{i-1}\oplus\rho_i\); repeated traces are therefore a
necessary part of a credible experimental model.

\paragraph{Interpretation.}
These experiments support the following claim and no stronger one.  The
uncompressed toric discrepancy has a real Montgomery/Kummer side-channel
instantiation: its exceptional divisor \(u=\pm1\) corresponds to the order-four
Curve25519 inputs that force the Libgcrypt ladder into low-order Kummer states,
and those states determine whether a specific field-arithmetic reduction is
short or long.  The real attack observed this reduction label with a cache
side channel and used the label sequence to recover scalar-bit transitions.  The
experiments above reproduce the algebraic identity, the deterministic ladder
label, and a timing realization of the same binary label.

The limitations are equally important.  First, this is an active chosen-input
ECDH setting, not the passive fixed-base secp256k1 signing setting.  Second, the
leakage label is binary and exceptional-divisor based; it is not one of the
compressed multi-class labels used elsewhere in the paper.  Third, the timing
surrogate demonstrates the label-to-bit mechanism but is not a re-exploitation
of a currently installed cryptographic library.  The historical implementation
claim rests on the published Libgcrypt attack and the CVE/advisory record,
whereas the supplied artifact provides a reproducible laboratory model of the
same discrepancy label.
